\par\Needspace{6\baselineskip}\medskip
\hypertarget{native-155-article111}{}\pdfbookmark[2]{111.}{article-111}
\begin{center}111.\end{center}
Wir nehmen nun zweitens an, die beiden Orte seien durch ihre Abstände von drei Ebenen gegeben, die einander in der Sonne rechtwinklig schneiden. Wir bezeichnen diese Abstände für den ersten Ort mit $x$, $y$ und $z$, für den zweiten mit $x'$, $y'$ und $z'$. Die dritte Ebene sei die Ekliptik selbst, und die positiven Pole der ersten und zweiten Ebene sollen bei den Längen $N$ und $90^0+N$ liegen. So gilt nach Artikel 53, wenn die beiden Radiusvektoren
\SourceLeaf{156}mit $r$ und $r'$ bezeichnet werden,
\[
\begin{aligned}
x&=r\cos u\cos(N-\AscendingNode)+r\sin u\sin(N-\AscendingNode)\cos i\\
y&=r\sin u\cos(N-\AscendingNode)\cos i-r\cos u\sin(N-\AscendingNode)\\
z&=r\sin u\sin i\\
x'&=r'\cos u'\cos(N-\AscendingNode)+r'\sin u'\sin(N-\AscendingNode)\cos i\\
y'&=r'\sin u'\cos(N-\AscendingNode)\cos i-r'\cos u'\sin(N-\AscendingNode)\\
z'&=r'\sin u'\sin i.
\end{aligned}
\]
Daraus folgt
\[
\begin{aligned}
zy'-yz'&=rr'\sin(u'-u)\sin(N-\AscendingNode)\sin i\\
xz'-zx'&=rr'\sin(u'-u)\cos(N-\AscendingNode)\sin i\\
xy'-yx'&=rr'\sin(u'-u)\cos i.
\end{aligned}
\]
Aus der Verbindung der ersten und zweiten Formel erhält man $N-\AscendingNode$ und $rr'\sin(u'-u)\sin i$. Daraus und aus der dritten Formel ergeben sich $i$ und $rr'\sin(u'-u)$.

Wird der Ort, dem die Koordinaten $x'$, $y'$ und $z'$ entsprechen, als der zeitlich spätere angenommen, so muss $u'$ größer sein als $u$. Steht außerdem fest, ob der um die Sonne vom ersten bis zum zweiten Ort überstrichene Winkel kleiner oder größer als zwei rechte Winkel ist, so müssen $rr'\sin(u'-u)\sin i$ und $rr'\sin(u'-u)$ im ersten Fall positiv und im zweiten negativ sein. Dann ist $N-\AscendingNode$ eindeutig bestimmt, und zugleich entscheidet das Vorzeichen von $xy'-yx'$, ob die Bewegung rechtläufig oder rückläufig ist. Ist umgekehrt die Bewegungsrichtung bekannt, so lässt sich aus dem Vorzeichen von $xy'-yx'$ entscheiden, ob $u'-u$ kleiner oder größer als $180^0$ zu nehmen ist. Sind hingegen sowohl die Bewegungsrichtung als auch die Beschaffenheit des um die Sonne überstrichenen Winkels völlig unbekannt, so kann man den aufsteigenden Knoten offenbar nicht vom absteigenden unterscheiden.

Man erkennt leicht, dass, ebenso wie $\cos i$ der Kosinus der Neigung der Bahnebene gegen die dritte Ebene ist, $\sin(N-\AscendingNode)\sin i$ und $\cos(N-\AscendingNode)\sin i$ die Kosinus der Neigungen gegen die erste beziehungsweise zweite Ebene sind. Ferner ist $rr'\sin(u'-u)$ der doppelte Flächeninhalt des von den beiden Radiusvektoren eingeschlossenen Dreiecks, und $zy'-yz'$, $xz'-zx'$ und $xy'-yx'$ sind die doppelten Flächeninhalte der Projektionen desselben Dreiecks auf die jeweiligen Ebenen.

Schließlich ist klar, dass anstelle der Ekliptik jede beliebige andere Ebene als dritte Ebene dienen kann, sofern alle durch ihre Beziehungen zur Ekliptik bestimmten Größen in gleicher Weise auf die jeweils gewählte dritte Ebene bezogen werden.
